\subsection{Strict morphisms of Fréchet spaces}

THEOREM 1. --- Let E and F be two Fréchet spaces and u a continuous linear mapping from E into F. The following conditions are equivalent :

\begin{enumerate}
\def\labelenumi{\alph{enumi})}
\item
  \(u\) is a strict morphism.
\item
  \(u\) is a strict morphism for the weakened topologies.
\item
  The image of u is closed in F.
\item
  \({}^t u\) is a strict morphism from \({\mathrm{F}}^{\prime }\) into \({\mathrm{E}}^{\prime }\) for the weak topologies.
\item
  The image of \({}^t u\) is closed in \({\mathrm{E}}^{\prime }\) for the weak topology \(\sigma \left( {{\mathrm{E}}^{\prime },\mathrm{E}}\right)\) .
\item
  The image of \({}^t u\) is closed in E' for the strong topology β(E', E).
\end{enumerate}

\(g)\) \({}^t u\) is a strict morphism from \(\mathrm{F}_c^\prime\) into \(\mathrm{E}_c^\prime\) (the duals endowed with the topology of compact convergence).

The equivalence of \(a\) , \(b\) ) and \(e\) ) follows from prop. 3 of IV, p.~28 and the remark preceding it. That of \(a\) ) and \(c\) ) is precisely cor. 3 of I, p.~19. The remark of IV, p.~28, also shows that \(d\) ) is equivalent to the fact that, the image of \(u\) is closed for the weakened topology \(\sigma(F, F')\) of \(F\) ; the equivalence of \(c\) ) and \(d\) ) then follows from prop. 2 of IV, p.~4.

We now prove the equivalence of \(e)\) and \(f)\) . It is enough to prove that \(f)\) implies \(e)\) . Suppose that the image of \({}^t u\) is closed for \(\beta(E', E)\) in \(E'\) . On account of the Banach-Dieudonné theorem (IV, p.~25, cor. 2), it is enough to prove that for every convex balanced neighbourhood \(U\) of \(0\) in \(E\) , the intersection \(B = {}^t u(F') \cap U\) is compact for \(\sigma(E', E)\) . The strong dual \(E_b'\) of the Fréchet space \(E\) is complete (IV, p.~22, prop. 2), hence the closed subset \(B\) of \(E_b'\) is complete, and so the normed space \(E_B'\) is complete (III, p.~8, corollary). Let \((V_n)\) be a decreasing sequence forming a fundamental system of neighbourhoods of \(0\) in \(F\) . Then \(F'\) is the union of sets \(C = V_n^\circ\) which are compact for \(\sigma(F', F)\) , hence \(E_B' = \bigcup B_n\) , with \(B_n = E_B' \cap {}^t u(C_n)\) . Since \(E_B'\) is a Baire space, and each of the sets \(B_{n}\) is convex balanced and closed, there exists a real number r \textgreater{} 0 and an integer n such that \(B \subset r \cdot B_{n}\) . Then we have \(\mathbf{B} = \mathbf{U}^{\circ} \cap {}^{t} u(r \cdot C_{n})\) ; since the sets \(U^{\circ}\) and \(r \cdot C_{n}\) are compact and \({}^{t} u\) is continuous for the weak topologies, B is compact for \(\sigma(\mathrm{E}', \mathrm{E})\) . This completes the proof of the equivalence of e) and f).

Finally the equivalence of \(g\) ) and the preceding conditions follows from prop. 18 of GT, X, § 2, No.~10 and the following lemma.

Lemma 1. --- Let E and F be two Hausdorff locally convex and quasi-complete spaces and u a continuous linear mapping from E into F. For \({}^t u\) to be a strict morphism from \(F_{c}'\) into \(E_{c}'\) , it is necessary and sufficient that the image \(u(E)\) of u be closed, and that every compact subset of \(u(E)\) be the image under u of a compact subset of E.

By Mackey's th. (IV, p.~2, th. 1) and the fact that on \(\mathrm{E}'\) (resp. F) the topology of compact convergence coincides with that of convex compact convergence (IV, p.~4), we can identify E (resp. F) with the dual \(\mathrm{E}_c'\) (resp. \(\mathrm{F}_c'\) ). Then \(u\) is the transpose of \({}^t u\) , and the equicontinuous subsets of E (resp. F) are the relatively compact sets. Lemma 1 then follows from prop. 2 (IV, p.~27), since \(u(\mathrm{E})\) is closed in F if and only if it is closed for the weakened topology \(\sigma(\mathrm{F}, \mathrm{F}')\) (IV, p.~4, prop. 2).

COROLLARY 1. --- Under the hypothesis of th. 1, the following conditions are equivalent :

\begin{enumerate}
\def\labelenumi{(\roman{enumi})}
\item
  \(u\) is a strict injective morphism;
\item
  \({}^{t}u\) is a strict surjective morphism for the weak topologies.
\item
\({}^{t}u\) is surjective.
\end{enumerate}

The implication (i) \(\Rightarrow\) (ii) follows immediately from the equivalence of conditions \(a\) , \(d\) ) and \(e\) ) of th. 1 and from IV, p.~6, prop. 5. It is clear that (ii) implies (iii). Finally, we prove that (iii) implies (i): if \({}^t u\) is surjective \(u\) is a strict morphism by the equivalence of \(a\) ) and \(e\) ) in th. 1; that \(u\) is injective follows from prop. 5 of IV, p.~6.

COROLLARY 2. --- Under the hypothesis of th. 1, the following conditions are equivalent :

\begin{enumerate}
\def\labelenumi{(\roman{enumi})}
\item
  \(u\) is surjective;
\item
  \(u\) is a strict surjective morphism;
\item
  \({}^t u\) is a strict injective morphism for the weak topologies.
\end{enumerate}

The equivalence of (i) and (ii) follows from Banach's th. (I, p.~17, th. 1).

In view of the equivalence of \(a\) and \(c\) in th. 1, condition (ii) says that \(u\) is a strict morphism and that its image is dense in F for \(\sigma(F, F')\) . The equivalence of (ii) and (iii) then follows from the equivalence of \(a\) and \(d\) in th. 1 and from prop. 5 of IV, p.~6.

If \(u: \mathrm{E} \to \mathrm{F}\) is a strict morphism of Fréchet spaces, the transpose \({}^t u\) is not necessarily a strict morphism from \(\mathrm{F}_b'\) into \(\mathrm{E}_b'\) (IV, p.~62, exerc. 3). However, we have the following partial result :

COROLLARY 3. --- Under the hypotheses of th. 1, the following property implies the properties a) to g):

\begin{enumerate}
\def\labelenumi{\alph{enumi})}
\setcounter{enumi}{7}
\tightlist
\item
  \({}^{t}u\) is a strict morphism from \(F_{b}^{\prime}\) into \(E_{b}^{\prime}\) .
\end{enumerate}

When E and F are both Banach spaces, or both Montel spaces, property \(h)\) is equivalent to the properties \(a)\) to \(g)\) of th. 1.

Suppose that \({}^t u\) is a strict morphism from \(\mathrm{F}_b'\) into \(\mathrm{E}_b'\) . We shall prove that the image H of \({}^t u\) is closed in \(\mathrm{E}_b'\) , from which the first assertion of cor. 3 will follow.

Let G be the closure of the image of u in F; the space G, with the topology induced by that of F assigned to it, is a Fréchet space. The mapping \(u: E \to F\) factorizes as \(u = j \circ v\) where j is the canonical injection from G into F and where \(v \in \mathcal{L}(E; G)\) . Then we have \({}^t u = {}^t v \circ {}^t j\) , where \({}^t j\) is surjective, by Hahn-Banach th. (II, p.~24, prop. 2); also, \({}^t v\) is injective since \(v(E)\) is dense in G (IV, p.~6, prop. 5). By hypothesis, the mapping \({}^t u\) from \(F_{b}'\) onto H is open; since \({}^t j\) is surjective and continuous, the mapping \({}^t v\) induces a homeomorphism from \(G_{b}'\) onto H. But the dual \(G_{b}'\) of the Fréchet space G is complete (IV, p.~22, prop. 2); consequently H is complete, hence closed in \(E_{b}'\) .

If E and F are Montel spaces, the strong topology on \(E'\) (resp. \(F'\) ) coincides with the topology of compact convergence, and h) is just a reformulation of g).

If E and F are Banach spaces, so are \(E_{b}^{\prime}\) and \(F_{b}^{\prime}\) , and condition h) is equivalent to f) by the equivalence of a) and c) applied to \(u: F_{b}^{\prime} \to E_{b}^{\prime}\) .

COROLLARY 4. --- Suppose E and F are Banach spaces. For \({}^t u\) to be surjective, it is necessary and sufficient that there exist a real number \(r > 0\) such that \(\| x\| \leqslant r.\| u(x)\|\) for all \(x\in E\) .

This is simply a reformulation of the equivalence of the conditions (i) and (iii) of cor. 1.

COROLLARY 5. --- Let E and F be two Fréchet spaces and u a continuous linear mapping

from E into F. The following conditions are equivalent :

\begin{enumerate}
\def\labelenumi{\alph{enumi})}
\item
  \(u\) is an isomorphism from E onto F.
\item
  u is an isomorphism from E onto F for the weakened topologies.
\item
  \({}^t u\) is an isomorphism from \({\mathrm{F}}^{\prime }\) onto \({\mathrm{E}}^{\prime }\) for the weak topologies.
\item
  \({}^{t}u\) is an isomorphism from \(F'\) onto \(E'\) for the strong topologies.
\item
  \({}^t u\) is an isomorphism from \({\mathrm{F}}_{c}^{\prime }\) onto \({\mathrm{E}}_{c}^{\prime }\) .
\end{enumerate}

Since an isomorphism is none other than a strict bijective morphism, the equivalence of \(a)\) and \(b)\) follows from the equivalence of conditions \(a)\) and \(b)\) of th. 1.

It is clear that a) implies each of the conditions c), d) and e).

Conversely, suppose that one of the conditions \(c\) ), \(d\) or \(e\) ) is satisfied. It follows from th. 1 and its cor. 3 that \(u\) is a strict morphism from E into F, and \({}^t u\) is evidently bijective. Let N (resp. I) be the kernel (resp. the image) of \(u\) . Since \({}^t u\) is bijective, we have \(\operatorname{Im} {}^t u = \operatorname{E}'\) and \(\operatorname{Ker} {}^t u = \{0\}\) , and so \(\mathrm{N}^\circ = \mathrm{E}'\) and \(\mathrm{I}^\circ = \{0\}\) by prop. 2 of IV, p.~27. But N (resp. I) is a closed vector subspace of E (resp. F), and the theorem of bipolars (II, p.~44) implies that \(\mathrm{N} = \{0\}\) and \(\mathrm{I} = \mathrm{F}\) , hence \(u\) is bijective. We have therefore proved that \(u\) is an isomorphism.

